\section{Methods}\label{sec:methods}

We consider on-policy reinforcement learning for autoregressive language models with
trajectory-level rewards. Let $\pi_\theta$ denote a language model parameterized by
$\theta$, generating a completion $y=(y_1,\dots,y_T)$ conditioned on a prompt $x$. After
sampling a full trajectory, the model receives a scalar reward $R(x,y)\in\mathbb{R}$
computed by an external verifier, such as exact-answer correctness or unit-test pass
rate. Our objective is
\begin{equation}
J(\theta)=\mathbb{E}_{x \sim \mathcal{D},\, y \sim \pi_\theta(\cdot \mid x)}[R(x,y)],
\end{equation}
where $\mathcal{D}$ denotes the prompt distribution. We focus on the regime most
relevant to recent reasoning-model training: sparse outcome rewards, on-policy sampling,
and no learned value function.

\subsection{Policy Gradient with Trajectory-Level Reward}

For a fixed prompt $x$, the standard REINFORCE estimator is
\begin{equation}
g_{\mathrm{rf}}(x,y)=R(x,y)\sum_{t=1}^{T}\nabla_\theta \log \pi_\theta(y_t \mid y_{<t},x).
\end{equation}
In practice, a baseline $b(x)$ is subtracted to reduce variance without changing the
expected gradient as long as $b(x)$ does not depend on the sampled action at token $t$:
\begin{equation}
g_{\mathrm{base}}(x,y)=\left(R(x,y)-b(x)\right)\sum_{t=1}^{T}\nabla_\theta \log \pi_\theta(y_t \mid y_{<t},x).
\end{equation}

We use prompt-level multi-sample baselines. Given $K$ sampled completions
$\{y^{(i)}\}_{i=1}^{K}$ for the same prompt $x$ with rewards $\{R^{(i)}\}_{i=1}^{K}$:
\begin{equation}
b_{\mathrm{RLOO}}^{(i)}(x)=\frac{1}{K-1}\sum_{j \neq i}R^{(j)}
\end{equation}
is the leave-one-out baseline used in RLOO-style estimators, while GRPO-style
normalization can be written as
\begin{equation}
A_{\mathrm{GRPO}}^{(i)}(x)=\frac{R^{(i)}-\mu_R(x)}{\sigma_R(x)+\varepsilon},
\end{equation}
where $\mu_R(x)$ and $\sigma_R(x)$ are the within-prompt mean and standard deviation of
the $K$ rewards. In both cases, the resulting scalar advantage is then broadcast
uniformly over all tokens in the sampled trajectory. This scalar-broadcast assumption is
precisely what we relax.

\subsection{Token-Level Credit Redistribution}

We introduce token-level weights $w_t(x,y)$ that redistribute a trajectory-level
advantage across tokens:
\begin{equation}
g_{w}(x,y)=A(x,y)\sum_{t=1}^{T} w_t(x,y)\,\nabla_\theta \log \pi_\theta(y_t \mid y_{<t},x),
\end{equation}
where $A(x,y)$ denotes either $R(x,y)-b_{\mathrm{RLOO}}(x)$ or $A_{\mathrm{GRPO}}(x,y)$
depending on the baseline choice.

The weights satisfy
\begin{equation}
\sum_{t=1}^{T}w_t(x,y)=1, \qquad w_t(x,y)\ge 0.
\end{equation}
This normalization keeps the total update magnitude comparable across weighting schemes
and isolates the effect of credit redistribution from simple rescaling.

Uniform token credit corresponds to $w_t=1/T$ for all $t$. Our goal is to replace this
uniform allocation with \emph{intrinsic} weights derived from quantities already
produced by the policy during generation. We do not train token-level reward models,
estimate prefix values, or build explicit search trees.

\subsection{Intrinsic Token-Weighting Mechanisms}

We define each method through an unnormalized salience score $\alpha_t(x,y)\ge 0$ and
then normalize within each trajectory:
\begin{equation}\label{eq:normalize}
w_t(x,y)=\frac{\alpha_t(x,y)+\varepsilon}{\sum_{k=1}^{T}\left(\alpha_k(x,y)+\varepsilon\right)},
\end{equation}
with a small $\varepsilon>0$ to avoid degenerate all-zero cases. In implementation, the
weights are treated as \emph{detached} scalars when multiplying the policy-gradient
term, so the update does not introduce second-order derivatives through the weighting
function itself. This keeps the optimization rule close in spirit and cost to standard
critic-free policy gradients.

\paragraph{Uniform Weighting (Baseline).}
\begin{equation}w_t^{\text{uniform}}=\frac{1}{T}.\end{equation}
Combined with an RLOO or GRPO baseline, this recovers the standard critic-free estimator
that assigns identical credit to every token in the sampled completion.

\paragraph{Surprisal Weighting.}
Our first intrinsic score is token surprisal,
\begin{equation}
\alpha_t^{\text{surp}}(x,y)=-\log \pi_\theta(y_t \mid y_{<t},x).
\end{equation}
This emphasizes low-probability decisions under the current policy. Intuitively, these
are tokens where the model commits to a less routine continuation and where uniform
credit assignment may be most diluted by predictable filler tokens.

\paragraph{Entropy-Reduction Weighting.}
Our second intrinsic score measures local entropy reduction. Let
\begin{equation}
H_t^{\mathrm{pre}}(x,y)=-\sum_{v}\pi_\theta(v \mid y_{<t},x)\log \pi_\theta(v \mid y_{<t},x)
\end{equation}
denote the predictive entropy before sampling token $y_t$, and let
\begin{equation}
\Delta H_t^{\mathrm{ent}}(x,y)=\max\left(0, H_t^{\mathrm{pre}}(x,y)-H_{t+1}^{\mathrm{pre}}(x,y)\right)
\end{equation}
for $t<T$. For the final token we set $\Delta H_T^{\mathrm{ent}}=0$. We then use
\begin{equation}
\alpha_t^{\text{ent}}(x,y)=\Delta H_t^{\mathrm{ent}}(x,y).
\end{equation}
This score highlights tokens after which the model's next-step distribution becomes
substantially more concentrated. Such events often correspond to commitment points in
reasoning, where one branch of continuation becomes much more likely than its
alternatives.

\paragraph{Policy-Divergence Weighting.}
A third intrinsic score measures how much the current policy has drifted from
a reference $\pi_{\mathrm{ref}}$ (typically the base model prior to RL) at
each token position:
\begin{equation}
\alpha_t^{\mathrm{div}}(x,y) = \lvert \log \pi_\theta(y_t \mid y_{<t}, x) -
\log \pi_{\mathrm{ref}}(y_t \mid y_{<t}, x) \rvert.
\end{equation}
In the RLHF/RLVR setting $\pi_{\mathrm{ref}}$ is already available because
it is used to compute the KL regularizer, so this score is essentially free.

\paragraph{Length-Robust Normalization.}
Because reasoning trajectories can vary substantially in length, all weights are
normalized within each sampled completion rather than across the minibatch. This ensures
that longer trajectories do not automatically receive larger aggregate updates simply
because they contain more token positions with nonzero salience. The comparison between
weighting schemes therefore reflects \emph{where} credit is assigned within a
trajectory, not how much total credit a longer sequence receives.

\paragraph{Batch-Level Baselines.}
All weighting schemes can be paired with either RLOO or GRPO-style prompt-level
baselines. Unless otherwise specified, we use the same rollout groups, optimizer,
sampling hyperparameters, and baseline family across methods. This isolates token
redistribution as the only algorithmic difference.

\subsection{Signal Collapse Under Low-Rank Adaptation}\label{sec:collapse}

So far our exposition has treated the policy $\pi_\theta$ as an unconstrained
language model being updated by gradient descent. In practice, however, almost
all LLM-RL pipelines apply LoRA~\citep{hu2021lora}, which restricts
$\pi_\theta$ to a rank-$r$ perturbation of a frozen reference policy
$\pi_{\mathrm{ref}}$. We now show that the intrinsic weighting schemes defined
above are qualitatively degraded by this constraint.

\paragraph{Notation.}
Write the logits produced by the base model at position $t$ as
$z_t^{\mathrm{ref}} = W_{\mathrm{lm}} h_t^{\mathrm{base}}$ and the logits
produced by the adapted model as $z_t^{\theta} = W_{\mathrm{lm}}
h_t^{\mathrm{adapted}}$, where $h_t^{\mathrm{adapted}} = h_t^{\mathrm{base}} +
\Delta h_t$ and $\Delta h_t$ is the sum of LoRA adapter contributions
propagated through the transformer to position $t$. Each adapter contributes
$B_\ell A_\ell x_{\ell,t}$ at layer $\ell$, where $B_\ell \in \mathbb{R}^{d
\times r}$ and $A_\ell \in \mathbb{R}^{r \times d}$, so the raw per-layer
perturbation at each position lies in the fixed $r$-dimensional column space
of $B_\ell$.

\paragraph{Collapse of surprisal and entropy.}
Because $\|\Delta h_t\|$ is bounded by a product of adapter norms and input
activation norms that are themselves roughly stationary across positions in a
well-trained base model, the first-order perturbation to per-token surprisal,
\[
-\log \pi_\theta(y_t \mid y_{<t}, x) =
-\log \pi_{\mathrm{ref}}(y_t \mid y_{<t}, x) + O(\|\Delta h_t\|),
\]
is dominated by the base model's surprisal pattern, which is task-agnostic
and primarily reflects filler versus content tokens. The same holds for
predictive entropy. Consequently, when we normalize these scores within a
trajectory (equation~\eqref{eq:normalize}), the resulting weights are nearly
identical to the base model's surprisal or entropy weights, regardless of
what the adapter has learned.

\paragraph{Collapse of policy divergence.}
Divergence weighting,
$\alpha_t^{\mathrm{div}} = \lvert \log \pi_\theta(y_t) - \log
\pi_{\mathrm{ref}}(y_t) \rvert$, is of direct interest because it quantifies
exactly the policy change that the RL update is trying to drive. Under LoRA,
however, this quantity is small and approximately uniform across positions for
a different reason: the KL budget $\mathrm{KL}(\pi_\theta \| \pi_{\mathrm{ref}})$
is bounded by the rank-$r$ adapter's capacity, so the per-token
log-probability ratios are compressed toward zero~\citep{anon2026implicitklra}.
Under the within-trajectory normalization in
equation~\eqref{eq:normalize}, this compression maps any approximately
constant score function to the uniform distribution, giving
$w_t^{\mathrm{div}} \approx 1/T$.

\paragraph{Consequence.}
We summarize the practical implication as follows. Let the Gini coefficient
$G(w)$ and the effective-token ratio $\mathrm{EffN}(w) / T = 1/(T \sum_t
w_t^2)$ be standard concentration measures, with $G(w) = 0$ and
$\mathrm{EffN}/T = 1$ corresponding to perfectly uniform weights. Then for
surprisal, entropy-reduction, and divergence weighting under LoRA, both
concentration measures approach their uniform limits as the adapter rank
$r$ decreases relative to the hidden dimension $d$. In this regime,
``intrinsic token weighting'' is indistinguishable from uniform broadcast,
and intrinsic-weighting methods cannot provide the credit-redistribution
benefit they were designed to deliver. This is \emph{not} a problem with the
underlying signals; it is a structural consequence of applying them inside a
low-rank adaptation.

\subsection{Adapter-Residual Credit Assignment (ARCA)}\label{sec:arca}

The collapse result motivates a different design choice. Rather than measuring
per-token properties of the \emph{output distribution} (whose shape is
dominated by the frozen base model under LoRA), we measure per-token
properties of the \emph{adapter's own contribution} to the forward pass.

\paragraph{Definition.}
Given a model with a LoRA adapter, let $h_t^{\mathrm{adapted}}$ denote the
last-layer hidden state at position $t$ with the adapter enabled, and
$h_t^{\mathrm{base}}$ the same hidden state with the adapter disabled.
Define the Adapter-Residual Credit Assignment salience as
\begin{equation}
\alpha_t^{\mathrm{ARCA}}(x, y) = \lVert h_t^{\mathrm{adapted}} -
h_t^{\mathrm{base}} \rVert_2,
\end{equation}
and normalize within the trajectory as in
equation~\eqref{eq:normalize} to obtain per-token weights
$w_t^{\mathrm{ARCA}}$. The adapter residual is computed once per minibatch
via an extra no-grad forward pass with the adapter disabled, using
\texttt{model.disable\_adapter()} in frameworks such as PEFT. This is exactly
the same call already required to compute a KL regularizer against the
reference policy or to support divergence weighting, so ARCA adds no new
infrastructure.

\paragraph{Why ARCA does not collapse.}
The key property that distinguishes ARCA from the intrinsic schemes above is
that it measures a quantity whose \emph{per-token} value is actively shaped by
the attention and MLP pathways through which the adapter input activations are
routed. Even for a low-rank adapter with small $\|B_\ell A_\ell\|$, the
input activations $x_{\ell,t}$ are non-uniform across positions (attention
patterns, layer norms, and content-versus-filler distinctions are already
non-uniform in the base model), so $\|\Delta h_t\|$ retains non-trivial
position-to-position variation. Formally, as the adapter rank goes to zero
the signal $\alpha_t^{\mathrm{ARCA}}$ vanishes uniformly, but its
\emph{normalized} counterpart $w_t^{\mathrm{ARCA}}$ remains non-degenerate
as long as the adapter is nontrivially active at \emph{any} position.

\paragraph{Interpretation as implicit gradient-informed credit.}
ARCA can be read as a cheap proxy for a gradient-based notion of token
importance. Positions at which the adapter contribution is large are
positions at which the adapter's parameter gradient has high contribution to
the loss, and therefore positions where a policy-gradient update can have
the most effect. This is a substantially weaker statement than the one made
by a learned critic, but it has the advantage of being available for free
from quantities the adapted forward pass already computes, and of being
exactly targeted at the PEFT-specific failure mode introduced above.

\subsection{Relationship to Existing Methods}

Our formulation recovers standard critic-free RL as a special case: uniform
weighting plus a prompt-level batch baseline yields the usual RLOO/GRPO-style
estimator. The intrinsic weighting schemes (surprisal, entropy reduction,
policy divergence) are the natural baselines within our framework and map
cleanly to published methods (e.g., divergence weighting corresponds to the
divergence-advantage analyses of \emph{Sparse but
Critical}~\cite{meng2026sparsebutcritical}, entropy-based weighting
to~\cite{he2026eapo,yu2026erpo}). ARCA is the only scheme we
study that uses an explicitly \emph{adaptation-aware} signal. Unlike
critic-based PPO, ARCA does not learn a value head. Unlike RED, PRIME, or
token-level reward-model methods, ARCA does not construct dense rewards from
an auxiliary model. Unlike TEMPO or related tree-based methods, ARCA does not
build prefix graphs or estimate branch values. Unlike hard token-selection
methods based on top-entropy masking, ARCA retains a dense token update but
modulates it continuously using an intrinsic salience score derived from the
adapter itself.

The resulting estimator is intentionally minimal. It changes only the
within-trajectory allocation of a scalar on-policy advantage, using
quantities already available from the forward pass with the adapter disabled.
This makes it easy to layer on top of existing critic-free RLVR pipelines
while preserving their computational simplicity.
